% Verification: Algorithms 4 (PipelineSameSkip), 5 (PipelineSameCross), and
% 6 (ApplyGamma) verified gate-by-gate against pipelined_gamma.py and
% Gamma.ipynb Construction 3.  All pass for L=3..15 (500 samples each).
% Depth formula: 8L+9 (odd L>=5), 8L+10 (even L>=6).
% Z correction condition: (L-1-c) mod 2 = 1, equivalently (t+L) is odd.

\subsection{Ancilla-free $\Gamma$ circuit}
\label{sec:gamma-impl}

We now construct an explicit circuit for $\Gamma$ using only nearest-neighbor $\CZ$, $\CNOT$, and $Z$ gates on the bare $L \times L$ grid, with \textbf{zero ancilla qubits}.
The circuit has depth $8L + O(1) = 8\sqrt{N} + O(1)$ and uses $O(N)$ gates---asymptotically negligible compared to the $O(N\sqrt{N})$ sorting stages.
(Empirically, the constant is $8L + 9$ for odd $L \ge 5$ and $8L + 10$ for even $L \ge 6$.)

We first describe the phase polynomial that $\Gamma$ must implement (\S\ref{sec:gamma-structure}), then develop two pipelined circuit constructions that realize it (\S\ref{sec:gamma-pipelines}), assemble the full circuit (\S\ref{sec:gamma-circuit}), and sketch the correctness argument (\S\ref{sec:gamma-correctness}).


\subsubsection{Phase polynomial and building blocks}
\label{sec:gamma-structure}

Recall that $\Gamma$ is diagonal: $\Gamma\ket{s} = (-1)^{f(s)}\ket{s}$.
We need $f(s)$ to be a degree-2 polynomial over $\mathrm{GF}(2)$ satisfying the parity-encoding condition of Lemma~\ref{lem:parity-encoding}.
The key observation is that $f(s)$ decomposes into a sum of \emph{triangular cross-products}, each involving data from at most two rows.
Define the primitive
\begin{equation}
T(x, y) \;=\; \bigoplus_{p < c'} x_p \cdot y_{c'},
\label{eq:primitive-T}
\end{equation}
where $x = (x_0, \ldots, x_{L-1})$ and $y = (y_0, \ldots, y_{L-1})$ are binary row vectors.
This is a strictly upper-triangular sum: every pair $(p, c')$ with $p < c'$ contributes the monomial $x_p \cdot y_{c'}$---a total of $\binom{L}{2}$ terms, far more than the $L{-}1$ nearest-neighbor connections on a single row.

The phase polynomial has two groups.
The first operates directly on qubit values:
\begin{equation}
f_D(s) = \bigoplus_{r \text{ even}} \Big[ \underbrace{T(s_r, s_r)}_{\text{same-row}} \oplus \underbrace{T(s_r, s_{r+1})}_{\text{cross-row}} \Big],
\label{eq:fD}
\end{equation}
pairing each even row $r$ with its odd neighbor $r{+}1$.
The second operates on \emph{column-parity} variables $\tilde{s}_{r,c} = \bigoplus_{r' \ge r} s_{r',c}$:
\begin{equation}
f_B(\tilde{s}) = \bigoplus_{\substack{r \text{ even} \\ r+2 \le L-1}} \underbrace{T(\tilde{s}_r, \tilde{s}_{r+2})}_{\text{skip-row}}
\;\oplus\; \bigoplus_{\substack{r \text{ even} \\ r \ge 2}} \underbrace{T(\tilde{s}_r, \tilde{s}_r)}_{\text{same-row}}.
\label{eq:fB}
\end{equation}

Three geometric configurations appear:

\paragraph{Same-row $T(x,x)$.}
All data on one row.
A prefix CNOT cascade builds running parities $\hat{x}_c = \bigoplus_{c' \le c} x_{c'}$.
A CZ between adjacent prefix-sum positions $\hat{x}_{c-1}$ and $\hat{x}_c = \hat{x}_{c-1} \oplus x_c$ contributes $\hat{x}_{c-1} \cdot \hat{x}_c = \hat{x}_{c-1} \cdot x_c + \hat{x}_{c-1}$ over $\mathrm{GF}(2)$.
Summing the degree-2 part over $c$ yields $\sum_{c} \hat{x}_{c-1} x_c = T(x,x)$; the degree-1 residual $\sum_{c=0}^{L-2} \hat{x}_c = \sum_p (L{-}1{-}p)\, x_p$ is corrected by single-qubit $Z$ gates on columns $p$ where $(L{-}1{-}p)$ is odd.

\paragraph{Cross-row $T(x,y)$.}
Row $r$ holds $x$; the adjacent row $r{\pm}1$ holds $y$.
A vertical CZ between prefix position $\hat{x}_c$ on row $r$ and untouched $y_c$ on row $r{\pm}1$ contributes $\hat{x}_c \cdot y_c = (\bigoplus_{c' \le c} x_{c'}) \cdot y_c$.
Summing gives $\bigoplus_{p \le c'} x_p\, y_{c'}$; a vertical CZ correction subtracts the diagonal $\bigoplus_c x_c y_c$, leaving $T(x,y)$.
Row $r{\pm}1$ is never modified.

\paragraph{Skip-row $T(x,y)$.}
Row $r$ (even) holds $x$; row $r{+}2$ (even) holds $y$; the odd row $r{+}1$ is a routing intermediary.
A 4-gate gadget per column---$\CZ(r{+}1, r{+}2)$, $\CNOT(r \to r{+}1)$, $\CZ(r{+}1, r{+}2)$, $\CNOT(r \to r{+}1)$---temporarily copies $\hat{x}_c$ onto the intermediary, applies the CZ interaction with row $r{+}2$, and restores the intermediary.
The two CZ gates cancel the intermediary's own state, isolating the desired $\hat{x}_c \cdot y_c$ contribution.
A matching correction gadget subtracts the diagonal terms.

\medskip
In all three cases, a prefix CNOT cascade on the source row builds the running parities $\hat{x}_c = \bigoplus_{c' \le c} x_{c'}$; the interaction gates extract the cross-column products; and an undo cascade restores the original data.
Executed alone, each would cost depth $2L + O(1)$.
The key to an efficient $\Gamma$ circuit is that these interactions can be \emph{pipelined}: the interaction gates trail the cascade wavefront at fixed column offsets, so multiple $T(\cdot,\cdot)$ terms sharing the same source row fold into a single forward-and-back sweep.


\subsubsection{Pipelined constructions}
\label{sec:gamma-pipelines}

The terms in $f_D$ and $f_B$ naturally group into two pipelined constructions, each performing a prefix cascade on an even row while interleaving trailing interaction gates that accumulate multiple $T(x,y)$ contributions simultaneously.

\paragraph{Construction~A: same-row + skip-row (Algorithm~\ref{alg:pipeline-a}).}
For each active even row $r \ge 2$ with $r{+}2 \le L{-}1$, this combines $T(\tilde{s}_r, \tilde{s}_r)$ and $T(\tilde{s}_r, \tilde{s}_{r+2})$ into one sweep.
(Row~$0$ uses a skip-only variant without the same-row CZ, since $f_B$ excludes same-row $T$ for $r = 0$.)
A forward prefix cascade on row $r$ advances left-to-right; at each time step $\tau$, multiple trailing operations execute in parallel at earlier (already-cascaded) columns.
The skip-row 4-gate gadget occupies offsets $-1$ through $-4$ behind the wavefront, routing the prefix through row $r{+}1$ to interact with row $r{+}2$.
The same-row CZ occupies offsets $-5$/$-6$, further back.
All trailing operations touch distinct columns at each step (Figure~\ref{fig:gamma-construction}(a)), so no qubit conflicts arise within a group.
The undo cascade sweeps right-to-left with mirrored trailing corrections.

Since the skip-row gadget accesses values on row $r{+}2$, groups whose target rows overlap cannot run simultaneously: group $r$'s gadget reads row $r{+}2$ while group $r{+}2$'s cascade would be modifying it.
We split into two batches: $r \equiv 0 \pmod{4}$ first, then $r \equiv 2 \pmod{4}$.
Within each batch, row triples are non-overlapping (e.g., $(0,1,2)$, $(4,5,6)$, $\ldots$\,).

\emph{Used in Phase~2 of Algorithm~\ref{alg:gamma}. Depth per batch: $2L + O(1)$. Two batches: $4L + O(1)$ total.}


\begin{algorithm}[t]
\caption{\textsc{PipelineSameSkip}: same-row $T(x,x)$ + skip-row $T(x,y)$ in one sweep}
\label{alg:pipeline-a}
\begin{algorithmic}[1]
\Input Even row $r$ (source), odd row $r{+}1$ (intermediary), even row $r{+}2$ (skip target).
  Data is in column-parity basis $\tilde{s}$.
\Output Applies phase $(-1)^{T(\tilde{s}_r,\,\tilde{s}_r) \;\oplus\; T(\tilde{s}_r,\,\tilde{s}_{r+2})}$.
  All rows restored.

\Statex
\Statex \emph{Forward sweep} --- at each time step $\tau$, the following
\Statex \hspace{1em} operations execute in parallel at distinct column offsets:

\Statex \hspace{2em}\textbf{Cascade} (offset 0): \hfill row $r$, cols $\tau{,}\,\tau{+}1$
\State \hspace{3em}$\CNOT\bigl((r,\,\tau) \to (r,\,\tau{+}1)\bigr)$
  \Comment{$q_{r,\tau+1} \gets \hat{x}_{\tau+1}$}

\Statex \hspace{2em}\textbf{Skip-row gadget} (offsets $-1$ to $-4$): \hfill rows $r{,}\,r{+}1{,}\,r{+}2$
\State \hspace{3em}$\CZ\bigl((r{+}1,\,\tau{-}1),\;(r{+}2,\,\tau{-}1)\bigr)$
  \Comment{pre-cancel}
\State \hspace{3em}$\CNOT\bigl((r,\,\tau{-}2) \to (r{+}1,\,\tau{-}2)\bigr)$
  \Comment{copy $\hat{x}$ onto intermediary}
\State \hspace{3em}$\CZ\bigl((r{+}1,\,\tau{-}3),\;(r{+}2,\,\tau{-}3)\bigr)$
  \Comment{interact: $\hat{x}_{\tau-3} \cdot y_{\tau-3}$}
\State \hspace{3em}$\CNOT\bigl((r,\,\tau{-}4) \to (r{+}1,\,\tau{-}4)\bigr)$
  \Comment{restore intermediary}

\Statex \hspace{2em}\textbf{Same-row CZ} (offsets $-5$/$-6$): \hfill row $r$
\State \hspace{3em}$\CZ\bigl((r,\,\tau{-}6),\;(r,\,\tau{-}5)\bigr)$
  \Comment{$\hat{x}_{\tau-6} \cdot x_{\tau-5}$}

\Statex
\State $\tau$ ranges from $0$ to $L{-}2$; out-of-bounds ops are skipped.
\State After $\tau = L{-}2$: drain trailing ops.
  \Comment{$O(1)$ steps}

\Statex
\Statex \emph{Undo sweep} --- reverse cascade (right-to-left) with trailing
\Statex \hspace{1em} corrections at mirrored offsets:

\Statex \hspace{2em}\textbf{Undo cascade:}
  $\CNOT\bigl((r,\,\tau{-}1) \to (r,\,\tau)\bigr)$
  \hfill restore $q_{r,\tau} \gets \tilde{s}_{r,\tau}$

\Statex \hspace{2em}\textbf{Skip-row correction gadget} (same 4-gate structure,
\Statex \hspace{3em} trailing at offsets $+1$ to $+4$)

\Statex \hspace{2em}\textbf{Same-row $Z$ correction} (offset $+5$):
  $Z(r,\,\tau{+}5)$ if $(L{-}1{-}(\tau{+}5))$ is odd
\end{algorithmic}
\end{algorithm}


\paragraph{Construction~B: same-row + cross-row (Algorithm~\ref{alg:pipeline-b}).}
For each even row $r$, this combines $T(s_r, s_r)$ and $T(s_r, s_{r+1})$ into one sweep.
The pipeline is shorter: the cross-row vertical CZ trails at offset $-2$ and the same-row CZ at offsets $-3$/$-4$ (Figure~\ref{fig:gamma-construction}(b)).
The undo sweep carries same-row $Z$ corrections and cross-row diagonal CZ corrections at mirrored offsets.

No batching is needed: cross-row CZ gates are diagonal and do not modify the computational-basis values on row $r{+}1$.
All even-row groups $(r, r{+}1)$ for $r = 0, 2, 4, \ldots$ are disjoint and execute in parallel.

\emph{Used in Phase~4 of Algorithm~\ref{alg:gamma}. Depth: $2L + O(1)$.}


\begin{algorithm}[t]
\caption{\textsc{PipelineSameCross}: same-row $T(x,x)$ + cross-row $T(x,y)$ in one sweep}
\label{alg:pipeline-b}
\begin{algorithmic}[1]
\Input Even row $r$ (source), odd row $r{+}1$ (cross-row target).
  Data is in original basis $s$.
\Output Applies phase $(-1)^{T(s_r,\,s_r) \;\oplus\; T(s_r,\,s_{r+1})}$.
  All rows restored.

\Statex
\Statex \emph{Forward sweep} --- at each time step $\tau$:

\Statex \hspace{2em}\textbf{Cascade} (offset 0): \hfill row $r$, cols $\tau{,}\,\tau{+}1$
\State \hspace{3em}$\CNOT\bigl((r,\,\tau) \to (r,\,\tau{+}1)\bigr)$
  \Comment{$q_{r,\tau+1} \gets \hat{x}_{\tau+1}$}

\Statex \hspace{2em}\textbf{Cross-row CZ} (offset $-2$): \hfill rows $r{,}\,r{+}1$
\State \hspace{3em}$\CZ\bigl((r,\,\tau{-}2),\;(r{+}1,\,\tau{-}2)\bigr)$
  \Comment{$\hat{x}_{\tau-2} \cdot y_{\tau-2}$}

\Statex \hspace{2em}\textbf{Same-row CZ} (offsets $-3$/$-4$): \hfill row $r$
\State \hspace{3em}$\CZ\bigl((r,\,\tau{-}4),\;(r,\,\tau{-}3)\bigr)$
  \Comment{$\hat{x}_{\tau-4} \cdot x_{\tau-3}$}

\Statex
\State $\tau$ ranges from $0$ to $L{-}2$; out-of-bounds ops are skipped.
\State After $\tau = L{-}2$: drain trailing ops.
  \Comment{$O(1)$ steps}

\Statex
\Statex \emph{Undo sweep} --- reverse cascade with trailing corrections:

\Statex \hspace{2em}\textbf{Undo cascade:}
  $\CNOT\bigl((r,\,\tau{-}1) \to (r,\,\tau)\bigr)$
  \hfill restore $q_{r,\tau} \gets s_{r,\tau}$

\Statex \hspace{2em}\textbf{Cross-row CZ correction} (offset $+2$):
  $\CZ\bigl((r,\,\tau{+}2),\;(r{+}1,\,\tau{+}2)\bigr)$

\Statex \hspace{2em}\textbf{Same-row $Z$ correction} (offset $+3$):
  $Z(r,\,\tau{+}3)$ if $(L{-}1{-}(\tau{+}3))$ is odd
\end{algorithmic}
\end{algorithm}


\subsubsection{Full $\Gamma$ circuit}
\label{sec:gamma-circuit}

Algorithm~\ref{alg:gamma} assembles the full $\Gamma$ circuit from four phases.
Phases~1 and~3 switch between the original and column-parity bases via vertical CNOT cascades.
Phase~2 applies $f_B$ using Construction~A (two batches); Phase~4 applies $f_D$ using Construction~B (single pass).

\begin{algorithm}[t]
\caption{\textsc{ApplyGamma}: ancilla-free $\Gamma$ on $L \times L$ grid}
\label{alg:gamma}
\begin{algorithmic}[1]

\Statex \textbf{Phase 1: Enter column-parity basis} \hfill Depth: $L{-}1$
\For{each column $c$ \textbf{in parallel}}
  \For{$r = L{-}2$ \textbf{down to} $0$}
    \State $\CNOT\bigl((r{+}1,c) \to (r,c)\bigr)$
  \EndFor
\EndFor
\Statex \hspace{1em} After this phase, qubit $(r,c)$ holds $\tilde{s}_{r,c} = \bigoplus_{r' \ge r} s_{r',c}$.

\Statex
\Statex \textbf{Phase 2: Parity-basis interactions} \hfill Depth: $4L{+}O(1)$
\Statex \hspace{1em}\emph{Batch 1:} for each even $r \equiv 0 \pmod{4}$ with $r{+}2 \le L{-}1$, \textbf{in parallel}:
\State \hspace{2em}\textbf{if} $r \ge 2$: \Call{PipelineSameSkip}{$r,\; r{+}1,\; r{+}2$}
  \Comment{same + skip}
\State \hspace{2em}\textbf{else} ($r = 0$): \Call{PipelineSkipOnly}{$r,\; r{+}1,\; r{+}2$}
  \Comment{skip only}
\Statex \hspace{1em}\emph{Batch 1 also:} for each even $r \ge 2$ with $r{+}2 > L{-}1$ not overlapping batch rows:
\State \hspace{2em}\Call{SameRowT}{$r$}
  \Comment{same-row only}
\Statex \hspace{1em}\emph{Batch 2:} for each even $r \equiv 2 \pmod{4}$ with $r{+}2 \le L{-}1$, \textbf{in parallel}:
\State \hspace{2em}\Call{PipelineSameSkip}{$r,\; r{+}1,\; r{+}2$}
  \Comment{Alg.~\ref{alg:pipeline-a}}

\Statex
\Statex \textbf{Phase 3: Exit column-parity basis} \hfill Depth: $L{-}1$
\For{each column $c$ \textbf{in parallel}}
  \For{$r = 0$ \textbf{to} $L{-}2$}
    \State $\CNOT\bigl((r{+}1,c) \to (r,c)\bigr)$
  \EndFor
\EndFor

\Statex
\Statex \textbf{Phase 4: Original-basis interactions} \hfill Depth: $2L{+}O(1)$
\Statex \hspace{1em}For each even $r$ with $r{+}1 < L$, \textbf{in parallel}:
\State \hspace{2em}\Call{PipelineSameCross}{$r,\; r{+}1$}
  \Comment{Alg.~\ref{alg:pipeline-b}}
\Statex \hspace{1em}If $L$ is odd: \Call{SameRowT}{$L{-}1$}
  \Comment{last even row, no cross partner}

\end{algorithmic}
\end{algorithm}


\paragraph{Conflict-freedom.}
In Phase~2, each batch is internally conflict-free: within a batch, row triples are non-overlapping.
Note that row~$0$ uses skip-row only (not \textsc{PipelineSameSkip}), since $f_B$ has same-row $T$ for even $r \ge 2$ only.
Remaining same-row-only rows ($r \ge 2$ with $r{+}2 > L{-}1$) are scheduled with the batch whose row triples they do not overlap.
The two batches execute sequentially to resolve the cross-group conflict on shared target rows.
In Phase~4, Construction~B groups use disjoint row pairs $(0,1)$, $(2,3)$, $\ldots$ and run fully in parallel.
If $L$ is odd, the last even row $L{-}1$ has no cross-row partner and uses a standalone same-row $T$ pass.

\paragraph{Resource summary.}
Total depth:
\[
\underbrace{(L{-}1)}_{\text{Phase 1}}
\;+\; \underbrace{2 \times (2L{+}O(1))}_{\text{Phase 2}}
\;+\; \underbrace{(L{-}1)}_{\text{Phase 3}}
\;+\; \underbrace{(2L{+}O(1))}_{\text{Phase 4}}
\;=\; 8L + O(1).
\]
Total gates: $O(N)$ ($L$ rows $\times$ $O(L)$ gates per row per phase).
Ancillas: \textbf{zero}.


\begin{figure*}[t]
    \centering
    \fbox{\parbox{0.95\textwidth}{\centering\vspace{3.0cm}
    \textbf{[PLACEHOLDER --- Composite figure]}\\[8pt]
    %
    \textbf{(a) Construction~A pipeline} (\textsc{PipelineSameSkip}, Algorithm~\ref{alg:pipeline-a}):\\
    Time ($\tau$) vs.\ column wavefront diagram for one row group $(r,\, r{+}1,\, r{+}2)$.
    Show cascade CNOT at offset $0$; skip-row 4-gate gadget at offsets $-1$ to $-4$
    (pre-cancel, copy, interact, restore on rows $r$, $r{+}1$, $r{+}2$);
    same-row CZ at offsets $-5$/$-6$.
    Mark which rows each operation touches.
    Undo sweep mirrors with correction operations.\\[8pt]
    %
    \textbf{(b) Construction~B pipeline} (\textsc{PipelineSameCross}, Algorithm~\ref{alg:pipeline-b}):\\
    Time vs.\ column wavefront diagram for one row pair $(r,\, r{+}1)$.
    Cascade at offset $0$;
    cross-row CZ at offset $-2$;
    same-row CZ at offsets $-3$/$-4$.
    Undo sweep with CZ and $Z$ corrections.\\[8pt]
    %
    \textbf{(c) Full $\Gamma$ overview}:\\
    Four-phase pipeline on the $L \times L$ grid.
    Phase~1/3: column CNOT cascades.
    Phase~2: two batches of Construction~A
    (annotate the cross-group conflict on row $r{+}2$ motivating batching).
    Phase~4: single pass of Construction~B.
    Annotate depth per phase.
    \vspace{3.0cm}}}
    \caption{The ancilla-free $\Gamma$ circuit.
    \textbf{(a,\,b)}~Pipelined constructions: a single prefix CNOT cascade on an even row absorbs multiple $T(x,y)$ interactions as trailing operations at fixed column offsets, achieving depth $2L{+}O(1)$ per sweep.
    In Construction~A, the skip-row 4-gate gadget routes interactions through an intermediary row; in Construction~B, cross-row interactions use direct vertical CZ gates.
    \textbf{(c)}~The four-phase structure of Algorithm~\ref{alg:gamma}. Phase~2 requires two batches because skip-row gadgets from group~$r$ access row $r{+}2$, which is the cascade row of group~$r{+}2$.
    Phase~4 needs no batching since cross-row CZ gates do not modify the target row.}
    \label{fig:gamma-construction}
\end{figure*}


\subsubsection{Why $\Gamma$ satisfies the parity-encoding condition}
\label{sec:gamma-correctness}

We sketch the argument for why Algorithm~\ref{alg:gamma} satisfies Lemma~\ref{lem:parity-encoding}; the complete monomial-level verification appears in Appendix~\ref{app:gamma-proof}.

The core question is: when we flip the two qubits involved in a vertical swap $(r_0, c_0) \leftrightarrow (r_0{+}1, c_0)$, does the phase polynomial change by exactly the parity string $P = \bigoplus_{\ell=j+1}^{k-1} s_\ell$?
The answer hinges on the \emph{parity} of $r_0$, which determines the shape of the JWT path between the two endpoints.

\paragraph{Even-row hop ($r_0$ even).}
The JWT path from $(r_0, c_0)$ to $(r_0{+}1, c_0)$ runs rightward along row $r_0$, turns the snake corner, and returns leftward along row $r_0{+}1$.
The parity string therefore spans columns $c' > c_0$ on both rows.
Because flipping both $s_{r_0, c_0}$ and $s_{r_0{+}1, c_0}$ cancels in the column-parity basis (the two flips offset), the parity-basis terms $f_B$ are unaffected: $\Delta f_B = 0$.
The entire contribution comes from $f_D$, specifically from $T(s_{r_0}, s_{r_0})$ and $T(s_{r_0}, s_{r_0{+}1})$.
In both terms, the flip at column $c_0$ toggles exactly those monomials involving columns $c' > c_0$---which is precisely the set of qubits in the parity string.

\paragraph{Odd-row hop ($r_0$ odd).}
The JWT path now runs leftward, so $P$ spans columns $c' < c_0$.
The column-parity variable $\tilde{s}_{r_0{+}1, c_0}$ does flip (since only one of the two qubits contributes at that level), and because $r_0{+}1$ is even, $f_B$ \emph{is} affected.
Six terms across $f_D$ and $f_B$ contribute to $\Delta f$.
All contributions from columns $c' > c_0$ cancel pairwise between the two groups, while contributions from columns $c' < c_0$ survive and sum to~$P$.
The parity-basis terms are precisely designed to supply the correction that the original-basis terms alone cannot provide for odd-row hops---this is the fundamental reason the construction needs both $f_D$ and $f_B$.